\section{The Pollux Framework}\label{sec:framework}

The first step is to define what it means for two descriptors to be compatible
with each other. This definition needs to be practical and flexible enough to
talk about real-world updates while also enabling formal reasoning about the
update process. To this end, Pollux considers two descriptors compatible if
there exists a well-behaved coercion between all messages of the writer's
descriptor into the reader's descriptor. The idea of a coercion as a function
also captures the idea that we should be able to reason about what message comes
out of the parser. If we don't know what message the parser might return, the
host program reading the encoding can't perform the necessary steps to recover
the original value, if that's desired. This framing is also grounded in
Programming Languages theory, were a line of research explored sub-typing is a
form of
coercion~\cite{breazu-tannenInheritanceImplicitCoercion1991,brandtCoinductiveAxiomatizationRecursive1998}.
Pollux extends this philiosphy, but rather than sub-typing between record types
(which our descriptors could easily be classified as), we're sub-typing through
the encoding format.

\subsection{Compatibility as Coercion}

Pollux's goal is to be able to look at a pair of data descriptors and
characterize what happens to values under the first descriptor when parsed under
the updated one. Our derivations, of the form
$\judg{\varnothing}{\dc_1 \dcompat \dc_2}$ provide more than a certificate of
compatibility between the descriptors: read rule by rule, it is a program
$\co : \denote{\dc_1} \rightarrow \denote{\dc_2}$, a coercion the reader applies to the
writer's message. We formalize this notion by making the coercion explicit, and
writing the judgment of the form
$\judg{\varnothing}{\dc_1 \dcompat \dc_2 \coerce \co}$.

But first, some definitions. Pollux operates on descriptors $\dc$, which are
written in a data descriptor language like Protobuf or CDDL. This can be thought
of as types for messages. On the other side is a message $\m$, which contains
some data while conforming to a descriptor. We'll use $\validv{m}{d}$ to
represent a message conforming to the descriptor, although the exact definition
of this will be format dependent. Likewise, $\denote{d}$ will be used as a type
to represent all messages which conform to the descriptor $\dc$.

I asserted earlier that two descriptors are compatible if a well-behaved
coercion exists. What does it mean for the coercion to be well-behaved? Clearly
we want to eliminate some coercions which arbitrarily change values in the
output message.

\vspace{1em}
\begin{definition}[Well-behaved Transformation]\label{def:well-behaved}
  A coercion $\co : \denote{\dc_1} \to \denote{\dc_2}$ is \emph{well-behaved} when,
  for all $\validv{\m}{\dc_1}$ and all $\bs$ with $\encodes{\dc_1}(\m, \bs)$,
  \[ \mathrm{parse}_{\dc_2}(\bs) = \some{(\co(\m), U)} \quad\text{with}\quad
    U \subseteq \dom{\dc_1} \setminus \dom{\dc_2}. \]
\end{definition}

Basically, a conforming descriptor and any possible encoding of that message
will yield the same message. Additionally, the parser reports which records in
the encoding were skipped (i.e. are writer only fields). The second condition
that $U \subseteq \dom{\dc_1} \setminus \dom{\dc_2}$ is essential to this
definition. Wire-soundness, that the message parses, would admit something like
$\ints \dcompat \String$ when the reader skips the mistyped record and injects
the default $\String$ value. This condition rules out the trivial transformation
which always injects the default values for all fields.

\vspace{1em}
\begin{definition}[Compatibility]\label{def:compat}
  Two descriptors $\dc_1$ and $\dc_2$ are compatible if there exists a well-behaved
  transformation between them. 
\end{definition}

There are two conditions part of Definition~\ref{def:well-behaved}, wire
soundness and faithfulness. Wire soundness states that a valid encoding of a
message under $\dc_1$ will parse to the coerced version under $\dc_2$. Faithfulness
states that there are no skipped fields during parsing.  An immediate
consequence of this definition is that coherence is free. Wire soundness pins
$f$ down on every valid message that has at least one encoding. So if
$\encodes{\dc_1}$ is total on valid messages (i.e. a serializer exists), any two
derivations of the same judgment yield coercions that agree on valid messages.

\subsubsection{Pollux Judgments}\label{sec:judgments}

A Pollux judgment has the form $\judg{\Gamma}{\nm_1 \dcompat \nm_2 \coerce \co}$ where
\begin{itemize}
\item $\Sigma$ is a symbol table, mapping names of descriptors to descriptor
  values. Representing this as a table allows us to store (mutually) recursive
  descriptors in finite syntax and mirrors Protobuf's own
  \texttt{FileDescriptorSet} construct.
\item $\Gamma$ is the assumption context, which stores pairs of descriptor names
  which are assumed to be compatible with each other. It maps pairs of
  descriptor names to \emph{assumption variables} which will hold the coercion
  between those descriptors in the completed derivation. An alternative
  formulation would be so that $\Gamma$ maps a pair of descriptors to a value of type
  $x : \denote{\Sigma(\nm_1)} \rightarrow \denote{\Sigma(\nm_2)}$ 
\item $\co$ is a function $\co : \denote{\Sigma(\nm_1)} \rightarrow \denote{\Sigma(\nm_2)}$, the coercion
  built by the derivation.
\end{itemize}

This is actually a stratified definition, yielding slightly different coercions
for each type of object:
\begin{align*}
  &\judg{\Gamma}{\nm_1 \dcompat \nm_2 \coerce \co}
  && \co : \denote{\Sigma(\nm_1)} \to \denote{\Sigma(\nm_2)} \text{ (message bodies)} \\
  &\judg{\Gamma}{\fd_1 \fieldrel \fd_2 \coerce \co}
  && \co : \denote{\fd_1} \to \denote{\fd_2} \text{ (fields: cardinality and type)} \\
  &\judg{\Gamma}{\ty_1 \typrel \ty_2 \coerce \co}
  && \co : \denote{\ty_1} \to \denote{\ty_2} \text{ (payloads)}
\end{align*}

The compatibility relation / coercions presented here aren't the only ones you
could define for Protobuf. Extra assumptions about how messages are encoded
(i.e. having a canonical encoding) may allow additional rules to be
admitted. Some rules are omitted not for causing wire-type incompatibilities but
because the results aren't meaningful, like converting from $\fixeds$ to
$\float$. Thus we claim that Pollux is sound, but not complete.

\subsection{Adding Recursion}\label{sec:recursive-types}

It is reasonable to ask why we add the indirection of a symbol table. Most
nested messages could be represented by directly embedding the descriptor into a
field that lives in the parent descriptor. This leads to a nice, natural
inductive structure which has strong reasoning principles. However, it makes it
impossible to represent recursive or mutually recursive Protobuf
messages. Indeed, since we can't model this type of Protobuf message we also
can't analyze any Protobuf messages which has a (mutually) recursive message in
it's transitive closure. In the data set, about 6.78\% of messages contain
recursion somewhere in their transitive closure (See
Table~\ref{tab:recur-summary} for more details). Additionally, there is the
question of how to get the Lean formalization to read the Protobuf
descriptors. There are no Lean bindings for Protobuf. While there are unofficial
C bindings for Protobuf, which might be compatible with the Lean FFI, getting
the correct representation of the descriptor into the formalization is
critical. Since Protobuf is a self-describing format, we can encode any Protobuf
descriptor as a Protobuf blob using an official Protobuf implementation like Go,
which can then be used with the Lean parser Pollux has developed. This provides
a higher level of trust. The issue is that \texttt{Descriptor.proto} is
self-recursive to handle nested messages.

Consider the addition of recursive descriptors to the inductive Pollux
formalization with nested descriptors embedded directly in the parent. We
immediately run into two problems. First, a recursive descriptor would require
an infinitely large descriptor to represent as we'd be forced to unroll the
recursive forever (all official Protobuf implementations only traverse 100
layers deep into a message, or 10,000 for the Go implementation for some reason,
so we could use this to bottom out the
recursion)\footnote{\url{https://protobuf.dev/programming-guides/proto-limits/}}
. With a descriptor rule like this one
\[
  \infer[D-Update]{ \forall\ \id \in \dom{\dc} \cap \dom{\dc'}.\ \dc[\id] \fieldrel \dc'[\id] }{ \dc \dcompat \dc' }
\]
(Note that for mandatory messages, which aren't supported in Protobuf, we'd need
that $\dom{\dc_{mand}} \supseteq \dom{\dc_{mand}'}$, i.e.\ that mandatory fields can
become optional.) the descriptor relation would become co-inductive since an
update to this descriptor would require showing $\dc_1[\id] \fieldrel \dc_2[\id]$ if
$\id$ is the field where we see the recursive type and the field relation would
include:
\[ \infer[F-Msg]{ \judg{\Gamma}{\dc_1 \dcompat \dc_2 \coerce \co} }{
    \judg{\Gamma}{\field{\card}{\dc_1} \fieldrel \field{\card}{\dc_2} \coerce \co
    }} \] So a derivation that $\dc_1[\id] \dcompat \dc_2[\id]$ would need to
include a derivation of itself! The natural solution here is to make $\dcompat$
co-inductive.

The flat symbol table representation avoids co-induction using two independent
finiteness facts:
\begin{enumerate}
\item Syntactically, a cyclic descriptor using a symbol table is representable
  in finite syntax. A symbol table stores \texttt{Stream}'s second field as the
  \emph{name} \texttt{Stream} rather than a copy of the descriptor. This is
  equivalent to a $\mu$-bound recursive type. What is infinite is what the
  descriptor unfolds to, and that unfolding is a regular tree. The checker
  exploit this by carrying an assumption set $\Gamma$ of pairs of message names where
  $(\nm_1, \nm_2) \in \Gamma$ records that
  $\Sigma(\nm_1) \dcompat \Sigma(\nm_2)$ has already been assumed. Every recursive call
  can add a pair, meaning $\Gamma$ never shrinks so $\Gamma$ can grow to at most
  $|\dom{\Sigma} \times \dom{\Sigma}| = \mathcal{O}(\dom{\Sigma}^2)$ before every recursive call would
  immediately succeed. Thus the derivation is finite.

  Thinking about this with the \texttt{Stream} example, checking \texttt{Stream}
  against \texttt{Stream'} takes two steps. In the initial check,
  $(\mathtt{Stream}, \mathtt{Stream'}) \not \in \Gamma$ so it is added and both
  descriptors are unfolded once. Field 1 gives $\ints \typrel \ints$ and field 2
  reduces to $\judg{\Gamma}{\mathtt{Stream} \dcompat \mathtt{Stream'}}$ which is now
  in $\Gamma$.
\item Semantically, all Protobuf values are finite even under a cyclic schema
  since they derive from an inductive grammar. However, two other facts bolster
  this claim in official Protobuf implementations. First, all nested messages
  can be omitted since they have presence checks and secondly, all official
  Protobuf implementations only traverse 100 layers deep into a message (or
  10,000 for the Go implementation for some
  reason)\footnote{\url{https://protobuf.dev/programming-guides/proto-limits/}}
  after which all parsing attempts will fail. Since our definitions always work
  with the idea that parsing will succeed, we can actually assume that the
  messages are finite. When the soundness proof re-encounters an assumed
  message, it has necessarily descended down at least one layer. Soundness can
  then use strong induction on the writer's message depth. We saw this with the
  \texttt{Stream} example, which reveals not an oddly named Protobuf message,
  but an important semantic fact we can take advantage of to argue for the
  finiteness of Protobuf messages.
\end{enumerate}

While I claimed this avoids co-induction, the symbol table is actually encoding
the co-inductive nature of the compatibility checker into a graph
traversal. Think of every descriptor in $\Sigma$ as a node, which are linked when we
determine two descriptors are compatible. The problem with embedding descriptors
directly into each other is when we have cycles in this graph. The symbol table
lets us shortcut walking a cycle forever. This aligns with in the co-inductive
notion of ``this property holds when there is no good reason for it not to''. If
every other field of the messages are compatible, then the recursive call will
be too.

\subsection{Other Properties for Pollux Compatibility Relations}\label{sec:compat-props}

Building a Pollux compatibility relation gives you the ability to check
compatibility between descriptors of the relevant data descriptor
language. However, there are some other properties which can provide beneficial
properties for your system.
\begin{itemize}
\item \textbf{Symmetry.} Our compatibility relations are focused on backwards
  compatibility, were an updated version of the system needs to be able to
  handle data generated by the old version. However, creating a symmetric
  compatibility relation and showing all your updates live inside it would
  provide backward and forward compatibility, were the old system would be able
  to interpret message from the new system was well. This might be particularly
  important during a distributed system rolling update. Forward compatibility
  can be checked by always considering downgrading a descriptor after every
  update, but that requires an additional derivation. 
\item \textbf{Transitivity.} A transitive compatibility relation would allow the
  system to be backwards compatible with \emph{any} previous version. A
  non-transitive compatibility relation would only provide backwards
  compatibility with the directly previous version. Removing fields can
  complicate transitivity, see Section~\ref{sec:desc-rel}.
\end{itemize}

\subsection{Data Descriptor Format Independent Rules}

This section details some general rules that should apply to most common data
descriptor languages.

\subsubsection{Names \& Recursion}

As discussed in Section~\ref{sec:judgments}, the context $\Gamma$ maps pairs of names
to assumption variables, which represent coercions (although since this is
building a co-inductive proof, we might use a variable to represent a coercion
we don't have a definition for yet).

\begin{mathpar}
  \infer[N-Assum]
    { \Gamma((\nm_1, \nm_2)) \mapsto x }
    { \judg{\Gamma}{\nm_1 \dcompat \nm_2 \coerce x} }

    \infer[N-Unfold]
    { (\nm_1, \nm_2) \notin \dom{\Gamma} \\
      \judg{\Gamma, (\nm_1, \nm_2) \mapsto x}{\Sigma(\nm_1) \dcompat \Sigma(\nm_2) \coerce \co} } {
      \judg{\Gamma}{\field{\kappa}{\nm_1} \dcompat \field{\kappa}{\nm_2} \coerce \mathsf{fix}\
        x.\ \co} }
\end{mathpar}

The \textsc{N-Assum} rules states that when we encounter two descriptor names
which are in the context $\Gamma$, we use the stored coercion between them. The
\textsc{N-Unfold} rule is what adds things to the $\Gamma$ context. Importantly, the
well-foundedness of the fix-point operator $\mathsf{fix}$ is dependent on the
fact that the names in the conclusion exist under some constructor which means
that $x$ is only every applied to nested messages of strictly smaller depth, in
this case burrowing the as-of-yet unintroduced notation for a Protobuf field.

\subsubsection{Descriptors}

Since Pollux can be used for other data descriptor languages as well, we wanted
to have a general rule which reasons about far-reaching updates to a descriptor. 

This rule is currently presented in a Protobuf-specific form, including an extra
condition that the \Oneof{} fields are correctly preserved. While Pollux tries
to reason about descriptor updates in a field-by-field manner, some language
features may require predicates which examine multiple fields at once, like
Protobuf's \Oneof{}.

\begin{mathpar}
  \infer[D-Update]
  { \forall \id \in \dom{\dc_1} \cap \dom{\dc_2}.\ \judg{\Gamma}{\dc_1[\id] \fieldrel \dc_2[\id] \coerce \co_{\id}} \\
    \oneofpres(\dc_1, \dc_2) }
  {
    \judg{\Gamma}{\dc_1 \dcompat \dc_2 \coerce \lambda \m.\
      \forall\ \id \in \dc_2.
      {\begin{cases}
        \id \mapsto \co_{\id}(\m[\id]) & \id \in \dom {\dc_1} \cap \dom{\dc_2} \\
        \id \mapsto \init(\dc_2[\id]) & \id \in \dom{\dc_2} \setminus \dom{\dc_1}
      \end{cases}}
  }
}
\end{mathpar}

The coercion between messages is complex to state, but simple to understand. If
each field at a key shared between the descriptors has a coercion, those
coercions are applied. For new keys in the reader, a default value is
injected. Keys that only the writer has are ignored, since the reader will skip
them using the wire type in the tag during parsing.

The OneofPreserved definition is part of Theorem~\ref{thm:trans-valid}. Placing
the one layer definition here means that a derivation
$\judg{\Gamma}{\dc_1 \dcompat \dc_2 \coerce \co}$ already contains OneofPreservedAll over
every pair that it actually touches during the derivation.

Once descriptors model the reserved set, a premise will need to be added to
\textsc{D-Update} requiring that the reserved set of $\dc_1$ is a subset of the
reserved set of $\dc_2$. % TODO

\textbf{Example.}

\begin{minted}{proto}
message Tree {
  int32 val = 1; Tree child = 2;
}
message Tree' {
  int64 val = 1; Tree' child = 2;
}
\end{minted}

\begin{enumerate}
\item \textsc{N-Unfold} on $(\mathtt{Tree}, \mathtt{Tree'})$ gives $\Gamma =
  (\mathtt{Tree}, \mathtt{Tree'}) \mapsto x$.
\item \textsc{D-Update} handles the shared keys 1 and 2.
  \begin{itemize}
  \item \textbf{Key 1.} \textsc{T-Varint} gives
    $\decop{\intl} \circ \encop{\ints}$, lifted by the IMP/IMP row
    of the cardinality table.
  \item \textbf{Key 2.} \textsc{F-Msg}, then \textsc{N-Assum}, gives $x$, lifted to the OPT slot as ``map $x$''.
  \end{itemize}
\item Result: $\mathsf{fix}\ x.\ \lambda \m.\ \{1 \mapsto \decop{\intl}\circ\encop{\ints}(\m[1]),\ 2 \mapsto \mathrm{map}\ x\ (\m[2])\}$
\end{enumerate}

This is the recursive function you would write by hand. $x$ is only ever applied
to \m[2], which is a nested message of strictly smaller depth. That's why the
note says the $\mathsf{fix}$ is well founded on the writer's value.

\subsubsection{Payloads}

\begin{mathpar}
  \infer[T-Refl]{ }{ \judg{\Gamma}{\ty \typrel \ty \coerce \mathrm{id}} }

  \infer[T-Varint]
    { \ty_1, \ty_2 \in \vints }
    { \judg{\Gamma}{\ty_1 \typrel \ty_2 \coerce \decop{\ty_2} \circ \encop{\ty_1}} }

  \infer[T-StrByt]{ }{ \judg{\Gamma}{\String \typrel \byt \coerce \mathrm{id}} }

  \infer[T-BytStr]{ \text{bytes are UTF-8} }{ \judg{\Gamma}{\byt \typrel
      \String \coerce \mathrm{id}} }

  \infer[F-Msg]
    { \judg{\Gamma}{\nm_1 \dcompat \nm_2 \coerce \co} }
    { \judg{\Gamma}{\mathtt{msg}\ \nm_1 \typrel \mathtt{msg}\ \nm_2 \coerce \co} }
\end{mathpar}

Pairs of wire-compatible but undesirable rules, such as $\double$ to $\fixedl$
get no rule.

Once again, there is no transitivity rule, which would allow us to arbitrarily
compose coercions. That is wire unsafe, for reasons explained
elsewhere. Transitivity can be a lemma about the plain relation (without
coercions) although that might not be useful. It cannot be a rule with a
coercion and thus can't be stated into the relation.

\subsection{Definitions \& Theorems}

To start with, we have to define compatibility for every level of our
compatibility relation. A type is semantically compatible with another if for
all values of that type, up to depth $n$, the coercion can perfectly translate
to the new type.

\vspace{1em}
\begin{definition}[Semantic compatibility, type level]\label{def:sem-compat-typ}
  A coercion $g : \denote{\ty_1} \rightarrow \denote{\ty_2}$ \emph{carries
    $\ty_1$ into $\ty_2$ up to depth $n$}, written
  $\semjudg[n]{g : \ty_1 \typrel \ty_2}$, when
  \[ \forall\ \validv{\pay}{\ty_1} \text{ with } \mathrm{depth}_f(\pay) \le n,\
    \forall\ \bs.\ \encodes{\ty_1}(\pay, \bs) \rightarrow
    \mathrm{parse}_{\ty_2}(\bs) = \some{g(\pay)}. \]
\end{definition}

\vspace{1em}
\begin{definition}[Semantic compatibility, field level]\label{def:sem-compat-fld}
  A coercion $g : \denote{\st_1} \rightarrow \denote{\st_2}$ \emph{carries $\st_1$ into $\st_2$ up
  to depth $n$}, written $\semjudg[n]{g : \st_1 \dcompat \st_2}$, when
\[ \forall\ \validv{\st}{\fd_1}, rs.\ \mathrm{depth}(\st) \le n \rightarrow \encodes[\id]{\fd_1}(\st, rs) \rightarrow
  \assemble_{\fd_2}^{\id}(rs) = \some{(g\ \st)} \]
and no record of $rs$ is skipped.
\end{definition}

Here $\encodes[\id]{\fd}(\st, rs)$ says that $rs$, a list of records with field
number $\id$, is a valid encoding of slot $\st$. The message-level relation is
defined from it:
\[
  \encodes{\dc}(\m, \bs) \iff
  \records(\bs) = \some{rs}
  \wedge \forall (\id, w, \pay) \in rs.\ \id \in \dom{\dc}
  \wedge \forall \id \in \dom{\dc}.\ \encodes[\id]{\dc[\id]}(\m[\id], \restr{rs}{\id}).
\]
Records of different fields may appear in any order. $\assemble_{\fd}(rs)$ is the
part of $\mathrm{parse}$ that builds one slot from that field's records
(Section~\ref{sec:philosophy}):
\begin{itemize}
\item An empty list gives $\init_{\ty}$, $\none$, or $[\,]$.
\item Singular slots combine the records' values with $\combine{\ty}$.
\item Repeated slots concatenate the records' values, packed or unpacked.
\end{itemize}
"No record of $rs$ is skipped" means
$\forall (\id, w, \pay) \in rs.\ \accepts(\fd_2, w)$.

\vspace{1em}
\begin{lemma}[Message from fields]\label{lem:msg-from-fields}
  Suppose $\mathrm{OneofPreserved}(\dc_1, \dc_2)$, and
  $\semjudg[n]{g_{\id} : \dc_1[\id] \fieldrel \dc_2[\id]}$ for every
  $\id \in \dom{\dc_1} \cap \dom{\dc_2}$. Let $f$ be the coercion
  \textsc{D-Update} builds from the $g_{\id}$. Then for all
  $\validv{\m}{\dc_1}$ with $\mathrm{depth}_v(\m) \le n$ and all $\bs$ with
  $\encodes{\dc_1}(\m, \bs)$,
  \[ \mathrm{parse}_{\dc_2}(\bs) = \some{(f(\m), U)} \quad\text{with}\quad
     U \subseteq \dom{\dc_1} \setminus \dom{\dc_2}. \]
\end{lemma}

The \textsc{F-Field} case of the inner induction reduces to one lemma per axiom,
none of which require the $\Sigma$ or $\Gamma$ context. If
$\forall\ j \le n. \semjudg[j]{\co : \ty_1 \typrel \ty_2}$ and
$\card_1 \cardrel{\ty_1}{\ty_2} \coerce L$ then
$\semjudg[n]{L(\co) : \field{\card_1}{\ty_1} \fieldrel \field{\card_2}{\ty_2}}$.

For message payloads, those lemmas would need $\co$ to be a merge homomorphism,
$\co(\mathrm{merge}\ \m\ \m') = \mathrm{merge}\ (\co\ \m)\ (\co\ \m')$. It's required when
updating from a repeated field to a singular one, i.e. \textsc{K-RepOpt} and
\textsc{K-RepImp}. It would also be required in \textsc{K-Opt} if Encodes may
split a singular message across several records (which is can, this is the merge
semantics of Protobuf). This property may come for free naturally, like
coherence.

There are also considerations here about if Encodes is closed under
concatenation. If some byte string $\bs$ encodes $\m$ and $\bs'$ encodes $\m'$, then
$\bs \doubleplus \bs'$ encodes $\mathsf{merge}(\m, \m')$. Then wire-soundness
applied to the concatenated byte string, together with $\mathsf{parse}$
respecting concatenation forces this equation on all valid messages. Validity of
$\mathsf{merge}(\m, \m')$ would need to check for oneofs, where the late member
wins.

\vspace{1em}
\begin{definition}[Semantic compatibility]\label{def:sem-compat}
  A coercion $\co : \denote{\dc_1} \rightarrow \denote{\dc_2}$ \emph{carries $\dc_1$ into
  $\dc_2$ up to depth $n$}, written $\semjudg[n]{\co : \dc_1 \dcompat \dc_2}$, when
  \[ \forall\ \validv{\m}{\dc_1} \text{ with } \mathrm{depth}_v(\m) \le n,\
    \forall\ \bs.\ \encodes{\dc_1}(\m, \bs) \rightarrow
    \mathrm{parse}_{\dc_2}(\bs) = \some{\co(\m)}. \]
  It \emph{carries $\dc_1$ into $\dc_2$}, written $\semjudg{f : \dc_1 \dcompat \dc_2}$, when
  it does so at every depth, and $\semjudg{\dc_1 \dcompat \dc_2}$ abbreviates
  $\exists\ f.\ \semjudg{f : \dc_1 \dcompat \dc_2}$.
\end{definition}

\vspace{1em}
\begin{definition}[Semantic reading of an Assumption Set]\label{def:sem-assum}
  $\semjudg[n]{\Gamma}$ holds when every assumed pair is semantically compatible at every
  strictly smaller depth:
  \[ \forall\ (\nm_1, \nm_2) \in \Gamma.\ \forall\ j < n.\ \semjudg[j]{ \coercion{\Sigma(\nm_1)}{\cdot}{\Sigma(\nm_2)} :
      \Sigma(\nm_1) \dcompat \Sigma(\nm_2)} \]
\end{definition}

\vspace{1em}
\begin{definition}[Soundness]\label{def:soundness}
  \[ \judg{\emptyset}{\dc_1 \dcompat \dc_2 \coerce \co} \implies \wfd{\dc_1} \implies \wfd{\dc_2}
    \implies \co \text{ is well behaved} \]
\end{definition} % TODO link this definition with semantic compatibility

\vspace{1em}
\begin{corollary}[Soundness of the compatibility relation]\label{cor:compat-sound}
  $\semjudg[n]{\emptyset}$ holds vacuously at every $n$, so
  \[ \judg{\emptyset}{\dc_1 \dcompat \dc_2} \implies\ \semjudg{ \coercion{\dc_1}{\cdot}{\dc_2} : \dc_1
      \dcompat \dc_2} \]
\end{corollary}


A proof of these theorems would require strong induction on $n$, since a nested
message might be of any depth less than the current message. At each inductive
step of depth $n$, a nested induction on the derivation of the descriptor
relation would be required.

\begin{itemize}
\item Induction on the $\judg{\Gamma}{\dc_1 \dcompat \dc_2}$ derivation itself is not enough, due to
  \textsc{N-Assum}. When doing induction on the derivation, the inductive
  hypothesis would state that the conclusion holds for all sub-derivations,
  which \textsc{N-Assum} doesn't have. We get the fact that $(\nm_1, \nm_2) \in \Gamma$ but
  not anything else to prove a semantic fact about the descriptors referenced by
  those names.
\item Induction on the depth of the message also doesn't work. The message depth
  is only decreased when we de-structure a nested message, but there are many
  different ways (like the \textsc{F-Msg} rule) which update sibling
  fields that doesn't decrease the depth. Without the depth decreasing, the
  inductive hypothesis can't be applied.
\end{itemize}
However both approaches used together would be able to discharge the proof, with
the outer induction on message depth able to handle when we step into a nested
message and the inner induction on the derivation able to walk all the sibling
fields being updated without a change in index.

\subsection{Not a Logical Relation}

In past version of Pollux attempted to state this as a logical relation. That
was the wrong framing for this project. A logical relation is a family of
relations defined via induction on types. Typically, a logical relation wants to
relate two different terms to the same semantic type. The semantic compatibility
definition I stated did not have those features, it was a flat statement that a
the encoded version of any message (of bounded depth) would parse to the
transformed value. The step indexing (a.k.a depth bounding) was just a strong
induction on the depth of the message and the ``fundamental theorem'' stated the
soundness of the syntactic judgment on descriptors against an independent
semantics capturing the coercion as already formalized above.

Start with the view that the descriptors are recursive record types (some
examples are discussed in Section~\ref{sec:recursive-types}).  Then the
descriptor relation I'm proposing would be a small proof system where we want to
show that two descriptors are compatible using the judgment
$\judg{\emptyset}{\dc_1 \dcompat \dc_2}$ (explained below), but we could \emph{define} what
it means for two descriptor to be compatible as
\begin{mathpar}
  \dc_1 := \mu\ X.\ F\ X

  \dc_2 := \mu\ X.\ G\ X

  \semjudg{\dc_1 \dcompat \dc_2 := \exists\ f: \denote{\dc_1} \rightarrow \denote{\dc_2}.\ \forall\
    \validv{\m}{\dc_1}, \bs.\ \encodes{\dc_1}(\m, \bs) \rightarrow \mathrm{parse}_{\dc_2}(\bs) =
    f(\m)}
\end{mathpar}

In English, then the definition of two descriptors being compatible is that
there exists a transformation such that for all valid messages under descriptor
$\dc_1$, we can serialize that message into some encoding that can be parsed to
produce the image of $\m$ under the transformation.

However, the lean formalization doesn't use the $\mu$-bound version of descriptor
types (See Section~\ref{sec:log-rel-lean}), but instead a flat symbol table
$\Sigma$. This symbol table is global and fixed by the \texttt{.proto} file,
containing names and descriptors for all the cycles the descriptor of interest
is part of. Additionally, there is $A$, a set of name pairs currently assumed to
be compatible; it grows as the derivation descends. The judgment is then
$\judg{\Gamma}{\dc_1 \dcompat \dc_2}$, with two rules for the assumption context $\Gamma$:
\begin{mathpar}
  \infer[N-Assum]{ (\nm_1, \nm_2) \in \Gamma }{ \judg{\Gamma}{\nm_1 \dcompat \nm_2 }}

  \infer[N-Unfold]
    { (\nm_1, \nm_2) \notin \dom{\Gamma} \\
      \judg{\Gamma, (\nm_1, \nm_2) \mapsto x}{\Sigma(\nm_1) \dcompat \Sigma(\nm_2) \coerce \co} }
    { \judg{\Gamma}{\nm_1 \dcompat \nm_2 \coerce \mathsf{fix}\ x.\ \co} }
\end{mathpar}

Compatibility of the top level descriptors is $\judg{\emptyset}{\dc_1 \dcompat \dc_2}$. The
\textsc{N-Unfold} rule is equivalent to $\mu$ unfolding, with a side condition
that $A$ grows.

Using that framing, we can define a logical relation as such:

\subsection{Formalizing in Lean}\label{sec:log-rel-lean}

Lean doesn't have co-inductive types, unlike Rocq, and only limited support in
newer versions for co-inductive predicates and no built-in support for things
like parameterized co-induction. While the native support for co-inductive
predicates might be enough to prove what Pollux cares about, we'd be close of
the boundaries of what lean supports, and if we exceed what the built-in
co-induction mechanisms support, it quickly devolves into having to hand craft
a greatest fix-point manually and state bisimulation principles manually. Rather
than that, I'd like to avoid co-induction entirely, which we've just argued is
possible.

The lean formalization could use a $\mu$-binder to actually model and interact
with recursive types, however this would be a lot of formalization and machinery
which isn't really needed. Instead, the lean formalization can use an additional
context $\Sigma$ which maps a name to a descriptor. All (mutually) recursive messages
are stored using their names in $\Sigma$ and fetched as needed. While $\Sigma$ itself
might be large for a large corpus, the number of mutually-recursive messages is
much smaller. The largest cycle observed in the data set is 22
mutually-recursive message descriptors, so $\mathcal{O}(|\Gamma|) = 22^2 = 484$

With the Lean formalizing using the symbol table approach, unfolding of the
descriptors corresponds to the explode operation (Section~\ref{sec:proto-desc})
and the induction rule for proving things on the descriptor is modeled with the
assumption context $\Gamma$.

The paper formalization can use the derivation of
$\judg{\Gamma}{\dc_1 \dcompat \dc_2 \coerce \co}$ to build the coercion, but lean can't
eliminate a \texttt{Prop} into a coercion. There are two alternatives here.
\begin{itemize}
\item Put the derivation into \texttt{Type}. Define \texttt{coerce : Deriv Σ Γ
    d₁ d₂ → Env Γ → Value → Value} by well-founded recursion on (value depth,
  derivation). This is a faithful rendering of the report, but it duplicates
  \texttt{Value.reinterpret}, and every lemma about \texttt{reinterpret} would
  have to be re-proved for \texttt{coerce}.
\item Keep the existing \texttt{Value.reinterpret}. Keep the relation as a
  \texttt{Prop} and prove agreement: a derivation implies that on valid
  messages, \texttt{Value.reinterpret d₁ d₂} equals the coercion the rules
  describe. In practice, the \texttt{getD} fallback at in the coercion function
  never fires on shared keys. That's the \texttt{Slot.reinterpret\_isSome} lemma
  already sketched in Section~\ref{sec:proto-transform}. The report presents the
  coercion rules as the specification, and Lean implements them with a function
  defined on the descriptors. Coherence is automatic because
  \texttt{reinterpret} doesn't depend on the derivation.
\end{itemize}

It's worth noting that since messages can always be omitted, having a cyclic
descriptor doesn't complicate message initialization,
Section~\ref{sec:msg-init}. Additionally, message validity
(Definition~\ref{def:vvalid}) isn't impacted by adding recursive messages since
it recurses on the values, which we've established are finite, using the explode
interface to access the descriptors as needed. Additionally, this means that
reinterpret function's termination metric should be switched to the writer's
value's depth and \texttt{Desc.OneofPreservedAll} becomes a finite conjunction
over reachable pairs which mirrors the structure of the relation itself. % TODO

Also, the lean parser combinator and theorem
\texttt{limitRecursiveStateCompat\_Correct} already handle most of the indexing
happening here.

% Local Variables:
% citar-bibliography: ("../pollux.bib")
% TeX-master: "../pollux.tex"
% TeX-command-default: "Make"
% jinx-local-words: "cardinality protobuf"
% End:
